\documentclass[11pt,a4paper]{article}
\usepackage{amsmath, amssymb, amsfonts}
\usepackage{graphicx}
\usepackage{geometry}
\usepackage{hyperref}
\geometry{margin=1in}

\title{Numerical Simulation of a 1D Blood Flow Test Problem}
\author{Your Name}
\date{\today}

\begin{document}

\maketitle

\begin{abstract}
This report presents a numerical implementation of a 1D blood flow test problem inspired by the manuscript ``Mathematical Modelling of 1D Blood Flow in Arteries with Aneurysm Using the Navier--Stokes Equations''. The test problem consists of a coupled system of hyperbolic PDEs modeling flow and cross-sectional area in an arterial segment. Numerical solutions are computed using the Lax--Wendroff finite difference scheme, and results are analyzed at selected time instances.
\end{abstract}

\section{Introduction}

The objective of this work is to implement a numerical solver for a simplified 1D blood flow system and analyze the propagation of initial localized pulses. The focus is on Section 5.3 of the manuscript. 

Key goals:

\begin{itemize}
    \item Implement finite difference methods (Lax--Friedrichs and Lax--Wendroff)
    \item Capture wave propagation
    \item Analyze damping and boundary conditions
\end{itemize}

\section{Mathematical Model}

The system of PDEs is:

\begin{equation}
Q_t + A_z + \frac{1}{5}Q = 0,
\end{equation}

\begin{equation}
A_t + Q_z = 0,
\end{equation}

where $Q(z,t)$ is the blood flow rate, $A(z,t)$ is the cross-sectional area, and $z \in (0,1)$.

\subsection{Initial Conditions}

The initial pulses are localized:

\begin{equation}
Q(z,0) =
\begin{cases}
> 0, & 0.4 - \epsilon < z < 0.4 + \epsilon \\
0, & \text{otherwise}
\end{cases}, \quad
A(z,0) =
\begin{cases}
> 0, & 0.7 - \epsilon < z < 0.7 + \epsilon \\
0, & \text{otherwise}
\end{cases},
\end{equation}

with $\epsilon = 0.02$.

\subsection{Boundary Conditions}

The boundary conditions are:

\begin{align}
Q_z(0,t) + A_z(0,t) &= 0, \\
Q_z(1,t) - A_z(1,t) &= 0,
\end{align}

implemented using finite differences.

\section{Numerical Method}

The system is solved using the **Lax--Wendroff finite difference scheme**, which is second-order accurate in both space and time.  

\begin{itemize}
    \item Spatial grid: $N = 600$ points, $z \in [0,1]$
    \item Time step: $dt = 0.0005$
    \item Output times: $t = 0.25, 0.5, 1.0, 1.7, 2.6$
\end{itemize}

The solver is implemented in \texttt{lax\_wendroff\_improved\_bc.py}. Plots are manually saved to \texttt{report/figures}.

Other schemes (Lax--Friedrichs, basic Lax--Wendroff) were implemented for comparison.

\section{Results}

The numerical results show the propagation of the initial pulses:

\begin{itemize}
    \item Flow pulse at $z = 0.4$ spreads in both directions.
    \item Area pulse at $z = 0.7$ propagates similarly.
    \item Damping term $\frac{1}{5}Q$ reduces the amplitude over time.
    \item Boundary conditions correctly reflect wave behavior at the domain edges.
\end{itemize}

\begin{figure}[h!]
\centering
\includegraphics[width=0.8\textwidth]{figures/Q_snapshots.png}
\caption{Evolution of $Q(z,t)$ at selected time instances.}
\end{figure}

\begin{figure}[h!]
\centering
\includegraphics[width=0.8\textwidth]{figures/A_snapshots.png}
\caption{Evolution of $A(z,t)$ at selected time instances.}
\end{figure}

\section{Discussion}

The system can be interpreted as two travelling waves:

\begin{equation}
Q + A \quad \text{and} \quad Q - A,
\end{equation}

propagating in opposite directions. The improved Lax--Wendroff scheme:

\begin{itemize}
    \item Reduces numerical diffusion
    \item Preserves pulse amplitude over time
    \item Enforces boundary conditions accurately
\end{itemize}

This makes it suitable for further experiments in blood flow modeling.

\section{Conclusion}

The implementation successfully:

\begin{itemize}
    \item Solves the coupled 1D blood flow PDE system
    \item Captures wave propagation and damping
    \item Produces reproducible snapshots for analysis
\end{itemize}

The repository is structured for future extensions, including new schemes, varying initial conditions, or additional source/boundary terms.

\end{document}